\subsection{R 的加法一致结构}

由于群 \(\mathbf{R}\) 是全序的，函数 \(x^{+}, x^{-}\) 与 \(|x|\) 在 \(\mathbf{R}\) 上的定义方式与在 \(\mathbf{Q}\) 上相同，并且满足上面就 \(\mathbf{Q}\) 列出的所有关系，特别是关系 (1) 至 (7)。设 \(a\) 是一个 \(>0\) 的实数，令 \(\mathrm{U}_{a}\) 为由所有对 \((x, y) \in \mathbf{R} \times \mathbf{R}\) （它们满足 \(|x - y| < a\) ）组成的集；当 \(a\) 取遍 \(>0\) 的实数组成的集（或数 \(1/n\) 组成的集）时，诸集 \(\mathrm{U}_{a}\) 组成实直线加法群 \(\mathbf{R}\) 的一致结构（称为实直线的加法一致结构）的一致邻域基本系。

函数 \(|x|\) 、\(x^{+}\) 与 \(x^{-}\) 在 R 上一致连续，而函数 sup \((x, y)\) 与 inf \((x, y)\) 在 \(R \times R\) 上一致连续；因此这些函数分别与把定义在 Q 与 \(Q \times Q\) 上的相应函数连续延拓所得的函数一致。

